\documentclass[../../../thesis.tex]{subfiles}
\begin{document}
\subsection{Obrázky}
V akademickej oblasti prírodných a~technických vied sa v~záverečných prácach
objavujú v~pomerne veľkom počte aj netextové grafické objekty. Patria sem grafy, schémy, diagramy, fotografie,
prípadne iné dvojrozmerné vizualizácie obsahu.
Nehovoríme o~ozdobných grafických prvkoch,
tie do práce tohto typu nepatria.
Obsahové grafické prvky budeme spoločne nazývať slovom obrázok.
Obrázok môže byť súčasť textového odseku,
ale tejto možnosti sa vyhýbame,
ak to nie je úplne nevyhnutné.
Uprednostňujeme tzv. plávajúcu formu obrázkov,
teda objektov, ktoré sa nemusia nachádzať bezprostredne
na mieste v texte, kde sú spomenuté.
V zdrojovom kóde umiestňujeme príkazy na sadzbu obrázku za odsek,
v~ktorom sa o~ňom hovorí po prvýkrát.
Pod každým obrázkom je textové označenie,
začínajúce skratkou slova obrázok (Obr.) a~nasleduje
poradové číslo obrázku v práci.

Na obrázku \ref{fig:measurement} je znázornený proces správneho merania výšky dieťaťa.
Grafický objekt je súčasťou plávajúceho prostredia \verb|figure|.
Samotnú grafiku pripravíme v externom editore,
exportujeme ju do niektorého z bežných formátov (JPG, PNG, PDF)
a~jej vloženie do finálneho PDF súboru
záverečnej práce zariadi makro
\verb|\includegraphics| z~balíka \verb|graphix|.
Automatické číslovanie má na starosti príkaz \verb|\caption|,
ktorého argument je text pod obrázkom.
\begin{verbatim}
\begin{figure}[!ht]
    \centering
    \includegraphics[scale=1.00]{img/Measurement.png}
    \caption{Pravidelné meranie výšky dieťaťa}
    \label{fig:measurement}
\end{figure}
\end{verbatim}

\begin{figure}[!ht]
    \centering
    \includegraphics[scale=1.00]{img/Measurement.png}
    \caption{Pravidelné meranie výšky dieťaťa}
    \label{fig:measurement}
\end{figure}

\subsubsection{Umiestnenie obrázkov}\label{sec:figPlacement}
Na obrázok sa v~texte odkazujeme prostredníctvom čísla.
Môžeme písať o~tom, že na obrázku 1 vidíme to a~to
alebo použijeme skratku -- obr. 1.
Slovo obrázok aj skratku píšeme v~odkaze v~texte
malým začiatočným písmenom, ak sa nachádza vo vnútri vety.
Na každý obrázok v~práci by mal existovať odkaz v~texte.

Umiestnenie obrázku v rámci dokumentu riadi pomerne
komplikovaný algoritmus, čo nie vždy vedie k uspokojivým výsledkom.
Polohu plávajúceho objektu môžeme čiastočne ovplyvniť
nepovinným parametrom prostredia \verb|figure|.
V príklade, ktorý sme uviedli si môžeme všimnúť prítomnosť parametra \verb|h!| v hranatých zátvorkách
na konci prvého riadka.
Predstavuje požiadavku, že preferujeme,  aby sa obrázok nachádzal vo výslednom dokumente presne na tomto mieste. \LaTeX\ niekedy umiestni obrázok na začiatok nasledujúcej strany, prípadne aj inam.
Ak sa mu obrázok nepodarí umiestniť, zaradí ho až na
samý koniec kapitoly aj spolu so všetkými nasledujúcimi obrázkami.
To nebýva žiaduce a žiaľ, nemáme príliš veľa možností, ako takýto výsledok ovplyvniť.
Pomôže zmena rozmerov obrázku, prípadne jeho premiestnenie inam v zdrojovom kóde.
Odporúča sa, aby sa prostredie \verb|\begin{figure}...\end{figure}| nachádzalo mimo textového odseku, t. j. treba ho od okolitého textu oddeliť minimálne jedným prázdnym riadkom zhora aj zdola.

Na ovládanie umiestnenia plávajúceho objektu môžeme použiť tieto
parametre \verb|h|~(here) -- umiestnenie v mieste výskytu,
\verb|t|~(top) -- umiestnenie na stránke hore, \verb|b|~(bottom) --
umiestnenie na stránke dole, \verb|p|~(page) -- umiestnenie
na samostatnej stránke na konci kapitoly. Prvé tri
parametre môžu byť doplnené znakom výkриčníka (\verb|!|),
ktorý požiadavku zosilňuje.
Jednotlivé parametre možno aj kombinovať.
Napríklad inštrukcia \verb|[!ht]| znamená, že chceme mať
obrázok na tom mieste, kde sa nachádza v zdrojovom kóde
a~ak to za žiadnu cenu nie je možné,
trebárs z~dôvodu nedostatku miesta pred koncom strany,
môže sa obrázok nachádzať aj v~hornej časti stránky.
Ani to však nemusí stačiť a~obrázok napokon nájdeme
na konci dokumentu.
V~takom prípade treba skúsiť niektoré
z~riešení spomenutých v~predchádzajúcom odseku.

\subsubsection{Označenie obrázku a text pod obrázkom}
Text pod obrázkom pozostáva z~označenia obrázku a z~vysvetľujúceho obsahu.
Mal by sa nachádzať spolu s~obrázkom na tej istej strane.

Text je dostatočne opisný,
aby bol jasný obsah obrázku aj pri rýchlom prechádzaní
práce bez nutnosti detailného čítania hlavného textu.
Ak opis pod obrázkom pozostáva iba z~jednej vety,
prípadne ide o~heslo bez vetnej štruktúry,
nepíšeme zaň bodku.
V~prípade viacerých viet už bodku alebo príslušné interpunkčné
znamienka použijeme na konci každej vety,
aj poslednej.
V príklade na obrázku \ref{fig:measurement} je text bez bodky
a~to je správne.

\subsubsection{Číslovanie a odkazy}
Obrázky číslujeme podľa výskytu v práci od čísla 1.
Používame jednoúrovňové číslovanie,
teda obrázok 1, obrázok 2, atď.
V~\LaTeX-u je automatické číslovanie obrázkov zabezpečené v definícii makra \verb|\caption|.

Odvolávanie sa na číslo obrázku rieši dvojica príkazov \verb|\label| a~\verb|\ref|.
Prvý príkaz zistí prítomnosť najbližšieho číselného
registra a priradí k nemu menovku, ktorú zadáme do argumentu.
Napríklad obrázok \ref{fig:measurement} má menovku \verb|fig:measurement|. Menovku volí autor textu, môže byť ľubovoľná, musí však začínať písmenom a nesmie obsahovať špeciálne znaky, ktoré majú v \TeX-u kategóriu vykonateľných príkazov (napr. \verb|\|, \verb|%|, \verb|#|, \verb|$|, zátvorky a podobne). Tiež treba venovať pozornosť tomu, aby sa rovnaká menovka nevyskytla v~texte v príkaze \verb|\label| viackrát, pretože by došlo k jej preťaženiu a znefunkčneniu odkazov.

Z praktických dôvodov sa ustálila prax začínať menovku skratkou typu číslovanej položky: \verb|fig| pre obrázok, \verb|eq| pri rovniciach, \verb|tab| ako menovka tabuľky, \verb|sec| v~prípade nadpisu, atď.
\end{document}